% Source: Wald GR, physical PDF p. 268
%         (printed p. 255).
% Coverage: Figure 10.2, the lapse function and shift vector.
% Figures/tables/footnotes: Figure 10.2; no tables or footnotes.
% Status: source-reviewed and compile-clean.
% Uncertainties: none.
\begingroup
\renewcommand{\thefigure}{10.2}
\begin{figure}[htbp]
\centering
\begin{tikzpicture}[scale=.92,>=Latex,font=\small]
  % Two neighboring hypersurfaces.
  \draw[thick] (-3.2,-1.35) .. controls (-1.8,-.75) and (1.5,-1.55) .. (3.2,-.95);
  \draw[thick] (-3.2,1.05) .. controls (-1.8,1.65) and (1.5,.85) .. (3.2,1.45);
  \node[left] at (-2.65,-1.02) {\(\Sigma_t\)};
  \node[left] at (-2.65,1.38) {\(\Sigma_{t+\Delta t}\)};

  \coordinate (O) at (0,-1.15);
  \coordinate (N) at (0,1.22);
  \coordinate (S) at (1.35,-1.18);
  \coordinate (T) at (1.35,1.20);
  \draw[->,thick] (O) -- (N) node[midway,left] {\(Nn^a\)};
  \draw[->,thick] (O) -- (S) node[midway,below] {\(N^a\)};
  \draw[->,thick] (O) -- (T) node[midway,right] {\(t^a\)};
  \draw[dotted,thick] (S) -- (T);
\end{tikzpicture}
\caption{A spacetime diagram illustrating the definition of the lapse function,
\(N\), and shift vector \(N^a\).}
\label{fig:10.2}
\end{figure}
\endgroup
